\documentclass[11pt,a4paper]{article}
\usepackage[margin=2.4cm]{geometry}
\usepackage{amsmath,amssymb}
\usepackage{booktabs}
\usepackage{longtable}
\usepackage{array}
\usepackage[table]{xcolor}
\usepackage{enumitem}
\usepackage{tcolorbox}
\usepackage[hidelinks]{hyperref}
\usepackage{titlesec}

\definecolor{slate}{HTML}{203040}
\definecolor{accent}{HTML}{176B91}
\definecolor{warn}{HTML}{B2560D}
\tcbset{
  notebox/.style={colback=accent!4,colframe=accent!60,boxrule=.6pt,
    left=6pt,right=6pt,top=5pt,bottom=5pt,fonttitle=\bfseries\small},
  warnbox/.style={colback=warn!5,colframe=warn!65,boxrule=.6pt,
    left=6pt,right=6pt,top=5pt,bottom=5pt,fonttitle=\bfseries\small}
}
\newcommand{\caution}[2]{\begin{tcolorbox}[warnbox,title={#1}]#2\end{tcolorbox}}
\newcommand{\notebx}[2]{\begin{tcolorbox}[notebox,title={#1}]#2\end{tcolorbox}}
\newcommand{\file}[1]{\texttt{\small #1}}
\titleformat{\section}{\large\bfseries\color{slate}}{\thesection}{.6em}{}
\titleformat{\subsection}{\normalsize\bfseries\color{slate}}{\thesubsection}{.6em}{}

\title{\textbf{Hurricane Helene and County Labor Markets}\\[.3em]
\large Data Construction and Estimation: A Complete Procedural Note}
\author{Piergiulio Fasciani}
\date{22 November 2026}

\begin{document}
\maketitle

\begin{abstract}
\noindent This note documents every procedure that stands between the raw
published records and the reported estimates: how each source is downloaded,
validated, filtered and merged; how the treated roster and the donor universe
are fixed; how the synthetic comparison is fitted; and the exact form of each
of the five regressions in the project. It is written to be read alongside the
code, and every rule stated here corresponds to a line in a named file. Where
a procedure is a judgement call rather than a mechanical step, the note says so
and gives the reasoning. Numerical values quoted are from the run of
20 November 2026.
\end{abstract}

\tableofcontents
\bigskip

%=====================================================================
\section{Scope and reading order}

The pipeline has two halves, and they are kept strictly apart. The first builds
a national county--month panel from raw published records and performs no
estimation whatsoever. The second takes that panel as given and estimates. No
stage in the first half ever sees an outcome regression, and no stage in the
second half ever modifies the panel.

\begin{center}
\small
\begin{tabular}{@{}lp{5.1cm}p{5.6cm}@{}}
\toprule
\textbf{Half} & \textbf{Stages} & \textbf{Produces} \\
\midrule
I. Construction & \file{src/preprocessing/01--07} & the $3{,}221\times57$ national panel \\
II. Estimation  & \file{src/labor/steps/01--07}, \file{src/migration/} & weights, five regressions, placebos \\
\bottomrule
\end{tabular}
\end{center}

\noindent The separation is enforced, not merely intended: \file{tests/test\_source\_first.py}
fails if an estimation artefact appears upstream, and stage 07 asserts the panel's
exact dimensions before saving it.

\notebx{The single most important convention}{%
\textbf{A missing value is never turned into a zero, and never interpolated.}
This recurs at every stage below. A county with no FEMA record is not a county
with zero verified loss; a month with no LAUS release is not a month of zero
employment. Each is held missing and the downstream filters decide what to do
with it. Where a source supplies a value the project declines to use, that
value is archived to an audit file rather than silently dropped.}

%=====================================================================
\part*{Part I --- Data construction and cleaning}
\addcontentsline{toc}{section}{Part I --- Data construction and cleaning}

\section{Sources}

\begin{center}
\small
\begin{tabular}{@{}llp{6.1cm}@{}}
\toprule
\textbf{Source} & \textbf{Vintage} & \textbf{Role} \\
\midrule
NOAA nClimGrid-Daily & Sept.\ 2024 & county precipitation, 25--30 Sept.\ 2024 \\
Census PEP           & 2024 & population denominator, county geography \\
Census ACS 5-year    & 2023 & households, median home value, median income \\
FEMA IHP valid registrations & v2, DR-4827--4832 & verified property and personal loss \\
BLS LAUS             & 2022m1--2026m9 & employment, labor force, unemployment \\
BLS QCEW             & 2022Q1--2026Q1 & workplace jobs, industry 10 \\
Census Vintage 2025  & 2020--2025 & annual population and migration components \\
\bottomrule
\end{tabular}
\end{center}

\caution{The LAUS route is not a direct BLS download}{%
The default ingestion route is \textbf{economicdata.dev}, an independent
distributor of BLS series, not \texttt{download.bls.gov}. This is recorded in
the build audit and is stated here because it is a dependency on a third party
for a headline input. An official-flat-file route exists in
\file{05\_labor.do} and is exercised by setting
\texttt{H\_LAUS\_SOURCE\_MODE = "official"}; it is dormant by default. Both
routes are validated against the same header contract before import.}

\section{Stage 01 --- NOAA precipitation}

Rainfall is the one exposure measure that contains no human reporting
behaviour, which is why it is built first and kept separate from FEMA.

\begin{enumerate}[leftmargin=1.3em,itemsep=2pt]
\item The monthly county file is downloaded, or a cached copy confirmed. If
      neither is available the stage exits with an error rather than
      proceeding on partial data.
\item Eight expected columns are confirmed by name before any parsing. A
      missing column aborts the stage.
\item County codes are stripped to digits and asserted to be exactly five
      characters that parse as a number.
\item \textbf{NOAA state prefixes are not Census FIPS.} NOAA uses its own
      state numbering, so an explicit eleven-row crosswalk maps the NOAA
      prefix to the Census state code; the county suffix is carried over and
      the resulting FIPS is treated as a \emph{candidate} until it is
      validated against the Census county name in stage 04.
\item Daily values of $-999$ or below are the source's missing sentinel and
      are set to missing.
\item The six-day total is formed, and then
      \textbf{set to missing unless all six daily values are present}. A
      partial week is not a smaller storm.
\end{enumerate}

\noindent Output: \file{noaa\_11states.dta}, one row per county, plus an audit
CSV of the candidate FIPS assignments.

\section{Stage 02 --- Census population and covariates}

Two independent pulls. The PEP file is filtered to \texttt{sumlev==50}
(counties), restricted to the eleven study states, and its state and county
codes are zero-padded to two and three characters before concatenation, which
prevents the classic loss of leading zeros.

The ACS pull goes through the Census API and is the one place where a silent
failure would be most damaging, because an API error returns a readable HTML
or JSON page that a CSV importer will happily parse into nonsense. The stage
therefore \textbf{opens the file and inspects its first line} for the expected
variable codes before importing it, and aborts if they are absent. Households,
median home value and median household income are destrung and any
non-positive value is set to missing.

\section{Stage 03 --- FEMA verified losses}

Six disaster declarations (DR-4827 through DR-4832) are queried at
registration level and aggregated to counties. Required columns are confirmed
per file. Real-property and personal-property verified loss are summed
separately, and crucially each is accompanied by a count of non-missing
contributing records; if a county contributed no observed value for a field,
the summed total for that field is \textbf{restored to missing} rather than
left at the zero that \texttt{collapse (sum)} would otherwise produce.

\caution{What FEMA verified loss is, and is not}{%
These are verified losses of households that \emph{applied} for Individuals
and Households Program assistance. They are not total economic storm damage,
and they embed application behaviour: propensity to apply varies with income,
insurance coverage, tenure and awareness. Any exposure measure built on them
inherits that selection. This is the motivation for reporting a rainfall-only
tercile split as a robustness check in Part II.}

\section{Stage 04 --- Reconciliation and the damage index}

\subsection{Geographic validation}

The NOAA file is merged to Census PEP. Two guards follow:

\begin{itemize}[leftmargin=1.3em,itemsep=2pt]
\item Any NOAA county that fails to match PEP is exported to an audit file and
      the stage \textbf{stops}. An unmatched county is a crosswalk failure, not
      a county to drop.
\item County \emph{names} are normalised (lower-cased, punctuation and the
      trailing word ``county'' removed) and compared. Disagreements are
      exported for review. This is the check that converts the candidate FIPS
      of stage 01 into a validated one: agreement of code \emph{and} name.
\end{itemize}

\subsection{The composite index}

For counties with both a FEMA record and complete rainfall, economic burden is
\[
\text{burden}_i \;=\; 100\times\frac{\text{verified loss}_i}{\text{households}_i\times\text{median home value}_i},
\]
the denominator being a housing-value proxy. Rainfall and burden are converted
to percentile rank scores on $[0,100]$ within the complete-case sample, and
\[
\mathrm{HDSI}_i \;=\; 0.4\times\text{rainfall score}_i \;+\; 0.6\times\text{burden score}_i .
\]

\caution{The 40/60 weighting is a judgement, and the index is descriptive}{%
The weights are not estimated and no procedure in this project selects them.
The index is used only to \emph{partition} already-treated counties into
terciles; it never defines treatment, and no estimate is a function of the
index's scale. The tercile split is therefore robust to any monotone
re-weighting that preserves the ordering, which is weaker than defending
$0.4/0.6$ itself. The rainfall-only split in Part II is reported precisely
because it removes both the weighting choice and the FEMA behaviour.}

\section{Stage 05 --- Labor data}

\subsection{LAUS}

Series are restricted to county records, parsed into the four measures
(unemployment rate, unemployed, employed, labor force), reshaped wide and
dated. The window is 2022m1--2026m9.

\subsection{QCEW}

Quarterly county files are filtered to the three codes that define total
covered employment: ownership \texttt{0}, industry \texttt{10}, aggregation
level \texttt{70}. The three monthly employment levels per quarter are kept and
reshaped to a monthly series.

\notebx{A deliberate defensive step}{%
After filtering, the import \textbf{whitelists only four variables} and
discards the thirty-plus other BLS fields, including \texttt{year} and
\texttt{quarter}. Those two are then regenerated from the loop counters. The
reason is that the generated period variables must be known-correct for every
quarter rather than inherited from a field whose formatting BLS has changed
across vintages.}

\subsection{The complete calendar}

A skeleton of every county crossed with all 57 months is constructed and the
observed data merged onto it, so that \textbf{missingness is explicit rather
than represented by an absent row}. The panel is asserted to be exactly
$57\times3{,}221 = 183{,}597$ rows; a change in national county geography
breaks the assertion rather than silently altering the donor universe.

\subsection{October 2025}

The CPS was not collected during the 2025 shutdown, so LAUS has no October 2025
release. Some mirrors nonetheless carry values for that month. The stage:

\begin{enumerate}[leftmargin=1.3em,itemsep=2pt]
\item flags every row where a source supplied an October 2025 LAUS value;
\item exports those rows to
      \file{audits/LAUS\_Oct2025\_source\_values\_excluded.csv};
\item sets all four LAUS measures to missing for that month and asserts it;
\item leaves QCEW untouched, since workplace jobs are an independent
      measurement unaffected by the CPS.
\end{enumerate}

\noindent October 2025 is consequently dropped from the estimation sample, and
the corresponding event-time indicator ($k=13$) is never created.

\section{Stage 06 --- Treatment and the donor universe}

\subsection{The frozen roster}

\begin{sloppypar}
The 39 treated counties are read from a dated reference file,
\file{treated\_NC\_county\_roster\_2026-11-13.csv}, asserted to contain
exactly 39 unique FIPS, and merged on.
\end{sloppypar} The roster is then \emph{compared} with the
FEMA-positive North Carolina definition regenerated from the current data; a
discrepancy stops the run when \texttt{STOP\_IF\_TREATMENT\_CHANGED} is set.

The point is to make treatment immune to later FEMA revisions and to any
reclassification that might occur after outcomes have been inspected. Treatment
is observational and exposure-defined, and freezing it is what prevents that
definition from drifting toward the results.

\subsection{Geographic exclusions}

Donors are excluded by state FIPS in three groups:

\begin{center}
\small
\begin{tabular}{@{}lp{9.2cm}@{}}
\toprule
\textbf{Group} & \textbf{States} \\
\midrule
Impact region & AL, FL, GA, IL, IN, KY, NC, OH, SC, TN, VA, WV \\
Buffer & AR, DC, IA, MD, MI, MS, MO, PA, WI \\
Other 2024 hurricanes & LA, TX \\
\bottomrule
\end{tabular}
\end{center}

\subsection{Coverage screening}

A donor must additionally have all 24 months of January 2022--December 2023
present and positive for employment, labor force and jobs, with a non-missing
unemployment rate.

\caution{Why the screen stops at December 2023}{%
The screen deliberately uses only the 24 training months. January--August 2024
is left untouched so that it remains a genuine holdout for the backdated audit
in Part II. Screening donors on 2024 coverage would contaminate that test.
Later outcomes never participate in donor selection.}

\noindent The pipeline requires at least 500 eligible donors and exactly 39
treated counties, and asserts both.

\section{Stage 07 --- The final panel}

Stage 07 adds event-time variables and re-asserts every invariant: dimensions,
FIPS length, 57 months per county, 39 treated at the reference month, 24
training months for every treated county, and that no county is both treated
and an eligible donor. It performs no estimation. The audit files it writes
record coverage by month and the full county-level treatment and donor
assignment.

%=====================================================================
\part*{Part II --- Estimation}
\addcontentsline{toc}{section}{Part II --- Estimation}

\section{The treated aggregate and the transformations}

The estimand is built on an \textbf{equal-county mean} of the 39 treated
counties, not a population-weighted state aggregate. The four fitting outcomes
are
\[
100\log(\text{employment}),\quad 100\log(\text{labor force}),\quad
100\log(\text{jobs}),\quad \text{unemployment rate (pp)} .
\]
Each county series is centred on its own 2022 mean. Each outcome is then scaled
by the treated aggregate's 2022 temporal standard deviation, subject to a floor
of $0.25$ for the three log outcomes and $0.15$ for the unemployment rate. The
realised scales were $0.915$, $0.804$, $1.644$ and $0.279$. The floors exist so
that an outcome which happened to be unusually flat in 2022 cannot dominate the
joint objective through a near-zero denominator.

\section{The synthetic comparison}

\subsection{Objective}

For a fitting window of $T$ months and $K=4$ outcomes, with $y_{kt}$ the
treated mean, $X_{ktj}$ donor $j$'s transformed outcome, $s_k$ the fixed scale
and $u_j = 1/J$ the uniform vector:
\[
\min_{w,\,a}\;
\frac{1}{2KT}\sum_{k,t}\left(\frac{y_{kt}-a_k-\sum_j X_{ktj}w_j}{s_k}\right)^{\!2}
\;+\;\frac{\lambda}{2}\sum_j\left(w_j-u_j\right)^2
\]
subject to $w_j\ge 0$ and $\sum_j w_j = 1$.

Two features deserve emphasis. The \textbf{outcome intercepts} $a_k$,
implemented by demeaning each outcome block over the fitting window, mean the
weights are fitted to pre-event \emph{movements} while a fixed treated/control
level difference is permitted. The ridge shrinks toward the \textbf{uniform}
vector rather than toward zero, which spreads weight across donors instead of
discarding them; the realised fit places positive weight on 115 donors with an
effective (inverse Herfindahl) count of $63.9$.

\subsection{Penalty selection}

$\lambda$ is chosen by fitting on 2022 and scoring on 2023 over the grid
$\{0.03,0.1,0.3,1,3,10\}$, ties resolving toward stronger regularisation.

\begin{center}
\small
\begin{tabular}{@{}lcccccc@{}}
\toprule
$\lambda$ & 0.03 & 0.1 & 0.3 & \textbf{1} & 3 & 10 \\
\midrule
2023 RMSE & .1543 & .1433 & .1236 & \textbf{.1094} & .1126 & .1198 \\
\bottomrule
\end{tabular}
\end{center}

\noindent $\lambda=1$ is selected. The definitive weights are then refitted
with intercepts over the full pre-event window, January 2022--August 2024.

\subsection{Solver}

The problem is solved in Mata by accelerated projected gradient (FISTA) with
Euclidean projection onto the simplex. The stopping rule is a genuine
optimality certificate rather than a step-size heuristic: for the simplex, the
duality gap
\[
\text{gap} \;=\; \max\!\left(0,\; g'w - \min_j g_j\right),
\qquad g = (A'A + \lambda I)w - c,
\]
bounds the distance from the optimum, and the solver stops when it falls below
$5\times10^{-5}$. The reported fit converged in 7{,}000 iterations with a gap
of $1.9\times10^{-5}$.

\subsection{What this is, and what it is not}

The method uses synthetic-DiD-style \textbf{unit} weights, outcome intercepts
and uniform pre-event time averaging. It does \textbf{not} estimate a separate
synthetic-DiD time-weight vector. Describing it as the full canonical
Arkhangelsky et al.\ estimator would overstate the implementation, and this
note does not.

\section{Validation of the comparison}

\subsection{The backdated out-of-sample audit}

Because the definitive weights are fitted over the entire pre-event window,
their pre-event fit is \textbf{in-sample and is not evidence}. The one genuinely
out-of-sample test refits the same objective with the same $\lambda$ on
2022--2023 only, and asks those weights to predict January--August 2024, which
they never saw.

\begin{center}
\small
\begin{tabular}{@{}lcccc@{}}
\toprule
& employment & labor force & jobs & unemp.\ rate \\
\midrule
Training RMSE (2022--23) & 0.108 & 0.098 & 0.365 & 0.050 \\
Forecast RMSE (2024m1--8) & 0.501 & 0.594 & 0.593 & 0.101 \\
Ratio & \textbf{4.66} & \textbf{6.08} & 1.63 & 2.00 \\
\bottomrule
\end{tabular}
\end{center}

\caution{This is the weakest result in the project and should be reported as such}{%
The pipeline's own caution threshold is a ratio of $2$, and employment and
labor force exceed it substantially. Two things should be said in mitigation,
and neither should be presented as a dismissal. First, the ratio is inflated by
a very tight training fit: in absolute terms a forecast error of $0.50$ log-\%
is roughly a quarter of the estimated aggregate immediate effect and about a
tenth of the high-tercile effect. Second, the degradation is informative in its
own right --- it is direct evidence that the comparison deteriorates as the
horizon lengthens, which is a reason to discount long-horizon estimates rather
than a reason to discount the immediate ones.}

\subsection{Pre-event tests}

Two tests are reported per outcome and they answer different questions.

\begin{itemize}[leftmargin=1.3em,itemsep=2pt]
\item \textbf{Drift.} A linear contrast across the pre-event coefficients,
      weighted by $k+16.5$ about the window midpoint. This is the test that
      bears on identification. It rejects for no outcome
      ($p \ge 0.84$ throughout).
\item \textbf{Strict joint zero.} A joint test that all 31 free pre-event
      coefficients are exactly zero. This rejects for every outcome, at
      month-to-month deviations of order $0.05$--$0.15$ units against
      estimated effects an order of magnitude larger. It is reported for
      completeness, not as a screen.
\end{itemize}

\noindent Five of six outcomes have zero individually significant pre-event
coefficients out of 32.

\subsection{Geographic placebos}

199 draws, seed 20261114. Each draw samples 39 eligible donors as a
pseudo-treated group, removes them from their own donor pool, and refits the
entire objective at the locked $\lambda$. The observed October--December 2024
shock is compared with the pre-event mean gap, as for the real estimate.

The observed shock is more extreme than \textbf{all 199 draws} for all four
fitting outcomes.

\caution{These are benchmarking ranks, not randomisation p-values}{%
Treatment was not randomly assigned across these pseudo-groups. The ranks say
that the observed movement is unusual relative to what geography and sampling
alone produce; they do not license a randomisation-inference interpretation.}

\section{The six regressions}

Every estimate in the project comes from one of six regressions. Each one
writes its own expanded specification --- command, regressors, sample, weights,
clustering, and the realised $N$, cluster count and residual degrees of freedom
--- to \file{outputs/audit/REGRESSIONS.txt} as it runs, so the record is
generated by the code and cannot drift from it.

\subsection{Shared structure}

The three labor regressions are the same model with the treated indicators
split differently:
\[
y_{it} \;=\; \alpha_i \;+\; \delta_t \;+\; \sum_{k\in\mathcal{K}} \beta_k\,
\mathbf{1}[i\in\mathcal{T}]\cdot\mathbf{1}[t = t^\ast + k] \;+\; \varepsilon_{it},
\]
estimated by \texttt{areg} with county fixed effects $\alpha_i$ absorbed,
calendar-month effects $\delta_t$, analytic weights, and standard errors
clustered on county.

\begin{center}
\small
\begin{tabular}{@{}lp{9.3cm}@{}}
\toprule
$t^\ast$ & September 2024 \\
$\mathcal{K}$ & $\{-32,\dots,18\}\setminus\{-1,0,13\}$, giving \textbf{48} indicators \\
weights & treated $1/39$; donor $\hat w_j$ from the fitted synthetic vector \\
sample & treated counties, plus donors with $\hat w_j > 10^{-10}$ \\
vce & \texttt{cluster(county\_id)}: 154 clusters, $df_r = 153$ \\
fitted & $N = 7{,}546$ \\
\bottomrule
\end{tabular}
\end{center}

\noindent Three event months carry no indicator, for three different reasons:
$k=-1$ (August 2024) is the omitted category that parameterises the model;
$k=0$ (September 2024) is the storm month and is dropped from the sample as a
transition month; $k=13$ (October 2025) has no LAUS data.

\subsection{The reference is not August 2024}

This is the single most misread feature of the specification. The omitted month
only parameterises the regression. Every reported coefficient is recentred to
the mean of all 32 pre-event months by a linear combination taken with its
joint covariance:
\[
\tilde\beta_k \;=\; \beta_k \;-\; \frac{1}{32}\sum_{m=-32}^{-1}\beta_m ,
\qquad \beta_{-1}\equiv 0 .
\]
A reported coefficient therefore answers ``relative to the typical pre-storm
gap'', not ``relative to August 2024''. The same recentring is applied
separately within each damage tercile.

\subsection{Regression 1 --- Main event study}

Six outcomes, one model each. Period effects average the recentred event
coefficients over three windows: immediate (2024m10--12), 2025 (2025m1--9) and
late matched (2025m11--2026m3). The two ratio outcomes, jobs/employment and
employment/labor force, are differences of log outcomes and are internally
consistent because all six share one donor weight vector.

These $6\times3=18$ estimates form the \textbf{main Holm family}.
\file{04\_joint\_event\_study.do}.

\subsection{Regression 2 --- Damage terciles}

Identical sample, weights, fixed effects, clustering and reference. The only
change is that each treated-by-event-month indicator is split three ways by
exposure tercile, 13 treated counties each, giving $3\times48=144$ regressors.

Because all three terciles sit in \textbf{one} regression, the high-minus-low
contrast is taken with the joint covariance rather than by differencing two
separately estimated models.

\begin{itemize}[leftmargin=1.3em,itemsep=2pt]
\item \textbf{Primary:} terciles of the composite HDSI.
\item \textbf{Robustness:} terciles of six-day rainfall alone, which contains
      no FEMA claim behaviour. Reported as a footnote comparison.
\end{itemize}

Each tercile effect is an ATT for that exposure group against never-treated
donors, so it requires only the parallel-trends assumption of Regression 1
applied within the group.

\caution{Exposure is not randomly assigned}{%
Damage correlates with terrain, income, industry mix and housing stock. The
high-minus-low contrast compares differently exposed counties; it is not a
randomised dose. The defensible claim is that the effect is graded in damage,
not that a marginal unit of damage causes a marginal unit of job loss.}

\file{04c\_dose\_event\_study.do}.

\subsection{Regression 3 --- Extended LAUS}

The same design carried to September 2026 on LAUS alone, past the point where
QCEW can match it. Five outcomes. Exploratory, and outside the main Holm
family. \file{04b\_extended\_laus.do}.

\subsection{Regressions 4 and 5 --- Migration}

A different panel: annual, July-to-July, county and year effects, clustered on
county.
\[
m_{cy} \;=\; \alpha_c \;+\; \delta_y \;+\; \sum_{y\neq 2024}\gamma_y\,
\mathbf{1}[c\in\mathcal{T}]\cdot\mathbf{1}[\text{year}=y] \;+\; u_{cy},
\qquad y \in \{2021,\dots,2025\}.
\]
2024 is the omitted last unexposed year and the reported effect is the 2025
coefficient; 2021--2023 serve as pre-tests. Regression 4 pools the 39 treated
counties; Regression 5 splits them by HDSI tercile in a single regression.

Two weighting schemes are estimated. The unweighted scheme compares eligible
counties equally; the synthetic scheme gives donors the fitted labor weights.
The \textbf{synthetic scheme is primary}, because the unweighted scheme fails
its pre-test for the two less-damaged terciles while the synthetic scheme
passes for all three.

\caution{Migration is an outcome, never a control}{%
Migration is realised after the storm. Conditioning on it would be a bad
control in the sense of Angrist and Pischke: it would absorb precisely the
adjustment channel the paper is about, and would bias the labor coefficients by
selecting on a post-treatment variable. It therefore appears only on the
left-hand side.}

\caution{A negative contrast is slower in-migration, not depopulation}{%
The treated counties are high-amenity net in-migration counties. Their net
domestic migration runs $12.3 \to 12.4 \to 12.0 \to 10.7 \to 8.2$ per 1{,}000
across 2021--2025, against roughly $0$--$5$ for the comparison group, and
population growth remains positive. A negative estimated contrast means these
counties attracted fewer new residents than they otherwise would have. They did
not depopulate, and the paper must not say they did.}

\file{src/migration/01\_census\_v2025.do}, \file{02\_migration\_dose.do}.

\subsection{Regression 6 --- Backdated-weight robustness}

The audit of Section 9.1 tests the fitting \emph{procedure}; it cannot test the
weights actually used, because those are fitted over the whole pre-event window
by construction. This regression closes that gap directly. It re-estimates both
Regression 1 and the tercile contrast of Regression 2 using the weight vector
fitted on \textbf{2022--2023 only}, which never saw January--August 2024.
Sample window, excluded months, fixed effects, clustering, recentring and the
three windows are unchanged; only the donor weights and the donor inclusion
rule differ.

The backdated weights are fitted on 24 months rather than 32 and are therefore
handicapped by construction, which makes this a conservative check.

\begin{center}
\small
\begin{tabular}{@{}llcccc@{}}
\toprule
& & \multicolumn{2}{c}{\textbf{Definitive (32m)}} & \multicolumn{2}{c}{\textbf{Backdated (24m)}} \\
\cmidrule(lr){3-4}\cmidrule(lr){5-6}
Outcome & Group & est. & $p$ & est. & $p$ \\
\midrule
\multicolumn{6}{@{}l}{\emph{Immediate window, 2024m10--12}}\\
Employment & low            & $-0.01$ & .992 & $-0.12$ & .889 \\
           & high           & $-4.27$ & .000 & $-4.38$ & .000 \\
           & high $-$ low   & $-4.26$ & .000 & $-4.26$ & .000 \\
Unemp.\ rate & high $-$ low & $+2.57$ & .000 & $+2.57$ & .000 \\
Emp/LF     & high $-$ low   & $-2.73$ & .000 & $-2.73$ & .000 \\
\multicolumn{6}{@{}l}{\emph{2025 window, 2025m1--9}}\\
Unemp.\ rate & high $-$ low & $+0.74$ & .000 & $+0.74$ & .000 \\
Emp/LF     & high $-$ low   & $-0.77$ & .000 & $-0.77$ & .000 \\
\bottomrule
\end{tabular}
\end{center}

\notebx{Why the contrasts barely move}{%
The high-minus-low contrast compares two treated groups against the
\emph{same} synthetic control, inside a single regression. Error in that
control enters both tercile coefficients and differences out. The contrasts are
therefore nearly invariant to which weight vector is used --- they agree to two
decimals --- while the individual tercile \emph{levels} do shift slightly. The
backdated audit's unflattering ratios bear on the aggregate ATT, not on the
damage gradient.}

\file{04d\_backdated\_weights\_robustness.do}.

\section{Inference}

\subsection{Clustering}

Standard errors are clustered on county throughout: 154 clusters in the labor
regressions, $df_r = 153$. Reported 95\% intervals use the $t$ critical value
at those degrees of freedom and are \emph{individual} intervals, not
simultaneous bands.

\subsection{Multiplicity}

Holm's step-down procedure is applied within each family. The main family is
the 18 outcome--window combinations of Regression 1; the tercile contrasts form
their own family; the extended-horizon estimates are outside both and are
labelled exploratory.

Raw and adjusted $p$-values are both retained in the output tables, so the
reader can apply whichever standard they prefer rather than being handed only
one.

\caution{The tests within a family are not independent}{%
Six outcomes computed from four underlying series cannot be independent: two of
them are exact identities of the others, so the correlation matrix of the six
has rank four. Holm is valid under arbitrary dependence, so it remains correct,
but it is conservative here --- it is controlling for 18 tests when the
effective number is smaller. This is an argument for reporting both raw and
adjusted values, which the pipeline does.}

\section{Reproducibility and audit}

\begin{itemize}[leftmargin=1.3em,itemsep=2pt]
\item \file{run\_all.do} builds everything from raw sources;
      \file{rerun\_estimation\_only.do} refits without rebuilding.
\item \file{outputs/audit/STAGE\_STATUS.csv} gates result families in the
      report. A file existing does not establish that its stage succeeded.
\item \file{outputs/audit/REGRESSIONS.txt} is regenerated on every run by the
      estimation stages themselves and records each fitted specification.
\item \file{SHA256SUMS.txt} pins the source files; raw-source hashes are kept
      separately.
\item Four offline checks run in CI on every push.
\item The readable report and chart atlas transform presentation only; they
      never recalculate weights, regressions or placebo fits.
\end{itemize}

\section{Scope of inference}

The design is observational and was revised after earlier outcomes had been
inspected. It is not preregistered, and fixed rosters and reproducible code do
not by themselves establish causal identification.

County clustering does not fully address correlated regional shocks or the
uncertainty introduced by estimating the donor weights, which is treated as
fixed at the inference stage. Geographic exclusion reduces but cannot eliminate
spillovers: commuting zones cross county and state lines. Source releases
revise, so the raw snapshots should be preserved for replication. Migration,
extended-horizon and damage-tercile estimates are exploratory relative to the
main family.

\end{document}
